\documentclass[12pt,a4paper]{article}
\usepackage{fontspec}
\usepackage{amsmath, amssymb, amstext}
\usepackage{graphicx}
\usepackage{geometry}
\usepackage{xcolor}
\usepackage{fancyhdr}
\geometry{margin=1in}
\pagestyle{fancy}
\setlength{\headheight}{15pt} 
\lhead{کوییز درس فیزیک مدرن}
\chead{دانشکده فیزیک دانشگاه صنعتی خواجه نصیرالدین طوسی}
\rhead{۱۴۰۵/۰۳/۰۱}
\usepackage{xepersian}
\settextfont{B Nazanin}

\begin{document}
	
	\begin{tabular}{|p{7cm}|p{7cm}|}
		\hline
		\textbf{نام و نام خانوادگی:} & \textbf{شماره دانشجویی:} \\
		\hline
	\end{tabular}
	
	
	\section*{سوالات:}
	\begin{enumerate}
		\item 
		دو دوقلو یک سفر رفت و برگشتی از زمین به ستاره‌ای که ۱۲ سال نوری فاصله دارد انجام می‌دهند. آلیس با سرعت \(0.6c\) سفر می‌کند. باب ۱۰ سال بعد از آلیس حرکت می‌کند و با سرعت \(0.8c\) سفر می‌کند.
		\begin{enumerate}
			\item[(الف)] نشان دهید که هر دو دوقلو در یک زمان به زمین بازمی‌گردند.
			\item[(ب)] کدام دوقلو هنگام بازگشت جوان‌تر است؟
		\end{enumerate}
		
		\item 
		سه چارچوب مرجع لخت \( S \) ، \( S' \) و \( S'' \) را در نظر بگیرید. فرض کنید \( S' \) با سرعت \( v \) نسبت به \( S \) و \( S'' \) با سرعت \( v' \) نسبت به \( S' \) حرکت می‌کند. تمام سرعت‌ها هم‌خط هستند.
		\begin{enumerate}
			\item[(الف)] معادلات تبدیل بین \( x, y, z, t \) و \( x', y', z', t' \) و همچنین بین \( x', y', z', t' \) و \( x'', y'', z'', t'' \) را بنویسید. این معادلات را با هم ترکیب کرده و روابط بین \( x, y, z, t \) و \( x'', y'', z'', t'' \) را به دست آورید.
			
			\item[(ب)] نشان دهید این روابط با معادلات تبدیل مستقیم \( S \) به \( S'' \) معادل هستند. سرعت \( S'' \) نسبت به \( S \) از قضیه جمع نسبیتی سرعت‌ها:
			\[
			v'' = \frac{v + v'}{1 + \dfrac{v v'}{c^2}}
			\]
			به دست می‌آید. این نشان می‌دهد که دو تبدیل متوالی لورنتس خود یک تبدیل لورنتس مستقیم است.
		\end{enumerate}
		\end{enumerate}
		
		\vspace{10pt}
		\textbf{موفق باشید.}
		
	\end{document}